\chapter{Related work}
\label{chap:ch2}

\par This chapter presents existing systems and research related to digital document handling and public administration workflows. The goal is to show how similar problems have been handled in other systems and what solutions already exist.

\par The systems and studies presented here focus on document submission, electronic workflows, and online public services. Some of them are real systems used by governments, while others are research papers that propose ways to improve administrative processes.

\par After presenting these systems, a short comparison is made to identify their limitations in relation to the problem addressed by this thesis. This helps explain the motivation behind Doclane and how it differs from existing solutions.

% -------------------------------------------------

\section{Automated electronic document management systems in municipal units}

Koptyakova et al. \cite{koptyakova2019} present a system for automated electronic document management in municipal institutions. The paper focuses on improving how public institutions handle documents and reducing the time needed for administrative work.

The authors explain that traditional document workflows in municipal units are slow and depend on manual work. Documents are usually registered, tracked, and processed by hand, which can cause delays when the number of requests is high.

To solve this, the paper proposes an automated system that digitizes document registration, tracking, and processing inside municipal institutions. This reduces manual work and makes the process faster.

The study also looks at the cost and benefits of such systems. It mentions that the payback period can be short because less manual work is needed and efficiency increases.

However, the system is mainly used inside one municipal institution. It focuses on internal document handling and not on the full process between citizens and institutions.

Compared to this, Doclane focuses on the full workflow between citizens and public institutions, including submission, validation, and feedback.

% -------------------------------------------------

\section{Digital document systems in public sector implementation}

Az Ab Aziz et al. \cite{abaziz2019ddms} study the implementation of a Digital Document Management System (DDMS) in the Malaysian public sector. The system is used to manage documents and records inside government institutions and aims to reduce paper use and improve document handling.

The paper shows that even when such systems exist, adoption is not always high. The main issues are not only technical, but also related to policy, training, security, and internal support.

The authors explain that successful implementation depends on clear guidelines and coordination between agencies. Without this, different institutions use the system in different ways, which reduces its effectiveness.

Compared to this, Doclane focuses more on workflows between citizens and institutions, but the same idea applies: a system is only useful if it is properly adopted and used in real work.

% -------------------------------------------------


\section{Document Management Systems in Public Sector Integration}

Electronic document management systems in public sector institutions are often part of larger information systems. International Journal of Advances in Engineering \& Technology \cite{dms_integration_2012} explains that government institutions use many systems such as finance systems, HR systems, CRM systems, and web portals. Because of this, document management systems need to work with other systems.

The authors explain that integration means systems can exchange data and work together, not that they are merged into one system. This is important because many public sector processes depend on data moving between systems.

The study also shows that integration can improve efficiency by reducing manual work and making workflows simpler. However, it is not easy. It can be expensive, take time, and may require changes in several systems. Some integration projects also fail due to poor planning.

Even with these problems, integration is still needed if institutions want to improve digital workflows.

Compared to this, Doclane also follows the idea of integration by connecting document workflows with different parts of an administrative system and improving communication between users and institutions.

% -------------------------------------------------

\section{Enter Finland}

Enter Finland (\url{https://enterfinland.fi}) is a government service used in Finland for immigration applications. Users can fill forms online, upload documents, and pay fees. After submission, some cases still require an in-person visit.

The system reduces the need for paper submission and physical visits at the start of the process. However, it is limited to immigration services and one type of institution.

Compared to this, Doclane is designed for different types of document workflows across multiple institutions.

% -------------------------------------------------

\section{DocuSign}

DocuSign (\url{https://www.docusign.com}) is a platform used for signing documents online. Users can upload documents, send them for signature, and track their status.

It removes the need to print and scan documents in many cases. It is mostly used in business environments.

However, it focuses mainly on signing and approval. It does not support full administrative workflows like structured requests or multi-step validation.

Compared to this, Doclane handles full workflows from submission to approval and feedback.

% -------------------------------------------------

\section{GOV.UK Services}

GOV.UK (\url{https://www.gov.uk}) is a UK government platform for online public services. Users can apply for documents, submit forms, and access different services.

It reduces the need for physical visits and allows many tasks to be done online. It covers many services in one platform.

However, each service has its own workflow and is not part of one unified system.

Compared to this, Doclane aims to provide a single workflow system for different types of document requests.

% -------------------------------------------------

\section{Comparison}

The systems described above solve parts of the same problem in different ways.

Enter Finland is focused on one domain. GOV.UK provides many services, but they are separate. DocuSign focuses on document signing.

All of them reduce paper use and improve digital access, but they are limited either by domain, structure, or function.

% -------------------------------------------------

\section{Gap}

There is no widely used system that handles document-based workflows across multiple public institutions in one unified system.

Most systems are either domain-specific, split into separate services, or focused on one function like signing.

Doclane addresses this by providing a unified workflow for submission, validation, feedback, and approval across institutions.